\Chapter{39}

\Verse{1}
{
    \zA{话说众人见平儿来了，都说：“你们奶奶做什么呢，怎么不来了？” 平儿笑道： “他那里得空儿来？ 因为说没得好生吃，又不得来，所以叫我来问还有没有，叫我
        再要几个拿了家去吃罢。” 湘云道：“有，多着呢！” 忙命人拿盒子装了十个极大的。 平儿道：“多拿几个团脐的。”
        众人又拉平儿坐，平儿不肯，李纨瞅着他笑道：“偏 叫你坐！” 因拉他身旁坐下，端了一杯酒，送到他嘴边。}
}
{
    \eA{Upon seeing, the story explains, P'ing Erh arrive, they unanimously
        inquired, "What is your mistress up to? How is it she hasn't come?" "How
        ever could she spare the time to get as far as here?" P'ing Erh smiled and
        replied. "But, she said, she hasn't anything good to eat, so she bade me, as
        she couldn't possibly run over, come and find out whether there be any more
        crabs or not;}
}
{
}

\Verse{2}
{
    \zA{平儿忙喝了一口，就要走， 李纨道：“偏不许你去!显见得你只有凤丫头，就不听我的话了。” 说着，又命嬷嬷 们：“先送了盒子去，就说我留下平儿了。”
        那婆子一时拿了盒子回来，说：“二奶 奶说：‘叫奶奶和姑娘们别笑话要嘴吃。 这个盒子里，方才舅太太那里送来的菱粉 糕和鸡油卷儿，给奶奶姑娘们吃的。’”
        又向平儿道：“说了：‘使唤你来，你就贪住 嘴不去了，叫你少喝钟儿罢。’” 平儿笑道：“多喝了，又把我怎么样？” 一面说， 一面只管喝，又吃螃蟹。}
}
{
    \eA{(if there be), she enjoined me to ask for a few to take to her to eat at
        home." "There are plenty!" Hsiang-yün rejoined; and directing, with
        alacrity, a servant to fetch a present box, she put in it ten of the largest
        crabs. "I'll take a few more of the female ones," P'ing Erh remarked. One
        and all then laid hands upon P'ing Erh and tried to drag her into a seat,
        but P'ing Erh would not accede to their importunities.}
}
{
}

\Verse{3}
{
    \zA{李纨揽着他笑道：“可惜这么个好体面模样儿，命却平常， 只落得屋里使唤。 不知道的人，谁不拿你当做奶奶太太看？” 平儿一面和宝钗湘云
        等吃喝着，一面回头笑道：“奶奶，别这么摸的我怪痒痒的。” 李氏道：“嗳哟!这硬 的是什么？” 平儿道：“是钥匙。”
        李氏道：“有什么要紧的东西怕人偷了去，这么 带在身上？ 我成日家和人说：有个唐僧取经，就有个白马来驮着他； 刘智远打天下， 就有个瓜精来送盔甲；
        有个凤丫头，就有个你。}
}
{
    \eA{"I insist upon your sitting down," Li Wan laughingly exclaimed, and as she
        kept pulling her about, and forcing her to sit next to her, she filled a cup
        of wine and put it to her lips. P'ing Erh hastily swallowed a sip and
        endeavoured immediately to beat a retreat. "I won't let you go," shouted Li
        Wan. "It's so evident that you've only got that woman Feng in your thoughts
        as you don't listen to any of my words!"}
}
{
}

\Verse{4}
{
    \zA{你就是你奶奶的一把总钥匙，还要 这钥匙做什么？” 平儿笑道：“奶奶吃了酒，又拿我来打趣着取笑儿了。” 宝钗笑道：“这倒是真话。
        我们没事评论起来，你们这几个，都是百个里头挑 不出一个来的。 妙在各人有各人的好处。” 李纨道：“大小都有个天理：比如老太太
        屋里，要没鸳鸯姑娘，如何使得？ 从太太起，那一个敢驳老太太的回？ 他现敢驳回， 偏老太太只听他一个人的话。 老太太的那些穿带的，别人不记得，他都记得。}
}
{
    \eA{Saying this, she went on to bid the nurses go ahead, and take the box over.
        "Tell her," she added, "that I've kept P'ing Erh here." A matron presently
        returned with a box. "Lady Secunda," she reported, "says that you, lady Chu,
        and our young mistresses must not make fun of her for having asked for
        something to eat;}
}
{
}

\Verse{5}
{
    \zA{要不 是他经管着，不知叫人诳骗了多少去呢!况且他心也公道，虽然这样，倒常替人上 好话儿，还倒不倚势欺人的。”
        惜春笑道：“老太太昨日还说呢，他比我们还强呢！” 平儿道：“那原是个好的，我们那里比得上他？” 宝玉道：“太太屋里的彩霞，是个 老实人。”
        探春道：“可不是‘老实’!心里可有数儿呢。 太太是那么佛爷似的，事 情上不留心，他都知道。 凡一应事，都是他提着太太行，连老爷在家出外去的一应
        大小事，他都知道，太太忘了，他背后告诉太太。”}
}
{
    \eA{and that in this box you'll find cakes made of water-lily powder, and rolls
        prepared with chicken fat, which your maternal aunt, on the other side, just
        sent for your ladyship and for you, young ladies, to taste. That she bids
        you," (the matron) continued, turning towards P'ing Erh, "come over on duty,
        but your mind is so set upon pleasure that you loiter behind and don't go
        back.}
}
{
}

\Verse{6}
{
    \zA{李纨道：“那也罢了。” 指着宝 玉道：“这一个小爷屋里，要不是袭人，你们度量到个什么田地？ 凤丫头就是个楚霸
        王，也得两只膀子好举千斤鼎，他不是这丫头，他就得这么周到了？” 平儿道：“先 时赔了四个丫头来，死的死，去的去，只剩下我一个孤鬼儿了。”
        李纨道：“你倒是 有造化的，凤丫头也是有造化的。 想当初你大爷在日，何曾也没两个人？ 你们看， 我还是那容不下人的？
        天天只是他们不如意，所以你大爷一没了，我趁着年轻都打 发了。}
}
{
    \eA{She advises you, however, not to have too many cups of wine." "Were I even
        to have too much," P'ing Erh smiled, "what could she do to me?" Uttering
        these words, she went on with her drink; after which she partook of some
        more crab. "What a pity it is," interposed Li Wan, caressing her, "that a
        girl with such good looks as you should have so ordinary a fortune as to
        simply fall into that room as a menial!}
}
{
}

\Verse{7}
{
    \zA{要是有一个好的守的住，我到底也有个膀臂了。” 说着不觉眼圈儿红了。 众人都道：“这又何必伤心，不如散了倒好。” 说着，便都洗了手，大家约着往
        贾母王夫人处问安。 众婆子丫头打扫亭子，收洗杯盘。 袭人便和平儿一同往前去。 袭人因让平儿到屋里坐坐，再喝碗茶去。 平儿回说：“不喝茶了，再来罢。”
        一面说， 一面便要出去。 袭人又叫住，问道：“这个月的月钱，连老太太、太太屋里还没放， 是为什么？”}
}
{
    \eA{But wouldn't any one, who is not acquainted with actual facts, take you for
        a lady and a mistress?" While she went on eating and drinking with
        Pao-ch'ai, Hsiang-yün and the other girls, P'ing Erh turned her head round.
        "Don't rub me like that!" she laughed, "It makes me feel quite ticklish."
        "Ai-yo!" shouted Li Wan. "What's this hard thing?" "It's a key," P'ing Erh
        answered.}
}
{
}

\Verse{8}
{
    \zA{平儿见问，忙转身至袭人跟前，又见无人，悄悄说道：“你快别问!横 竖再迟两天就放了。” 袭人笑道：“这是为什么，唬的你这个样儿？” 平儿悄声告诉
        他道：“这个月的月钱，我们奶奶早已支了，放给人使呢。 等别处利钱收了来，凑 齐了才放呢。 因为是你，我才告诉你，可不许告诉一个人去！”
        袭人笑道：“他难道 还短钱使？ 还没个足厌？ 何苦还操这心？” 平儿笑道：“何曾不是呢。 他这几年，只 拿着这一项银子翻出有几百来了。}
}
{
    \eA{"What fine things have you got that the fear lest people should take it
        away, prompts you to carry this about you? I keep on, just for a laugh,
        telling people the whole day long that when the bonze T'ang was fetching the
        canons, a white horse came and carried him!}
}
{
}

\Verse{9}
{
    \zA{他的公费月例又使不着，十两八两零碎攒了，又 放出去，单他这体己利钱，一年不到，上千的银子呢。” 袭人笑道：“拿着我们的钱，
        你们主子奴才赚利钱，哄的我们呆等着！” 平儿道：“你又说没良心的话，你难道还 少钱？” 袭人道：“我虽不少，只是我也没处儿使去，就只预备我们那一个。”
        平儿 道：“你倘若有紧要事用银钱使时，我那里还有几两银子，你先拿来使，明日我扣 下你的就是了。” 袭人道：“此时也用不着。}
}
{
    \eA{That when Liu Chih-yüan was attacking the empire, a melon-spirit appeared
        and brought him a coat of mail, and that in the same way, where our vixen
        Feng is, there you are to be found! You are your mistress' general key; and
        what do you want this other key for?" "You've primed yourself with wine, my
        lady," P'ing Erh smiled, "and here you once more chaff me and make a
        laughing-stock of me."}
}
{
}

\Verse{10}
{
    \zA{怕一时要用起来不够了，我打发人去取 就是了。” 平儿答应着，一径出了园门，只见凤姐那边打发人来找平儿，说：“奶奶有事 等你。”
        平儿道：“有什么事这么要紧？ 我叫大奶奶拉扯住说话儿，我又没逃了，这 么连三接四的叫人来找！” 那丫头说道：“这又不是我的主意，姑娘这话自己和奶奶
        说去。” 平儿啐道：“好了，你们越发上脸了！” 说着走来。 只见凤姐儿不在屋里，
        忽见上回来打抽丰的刘老老和板儿来了，坐在那边屋里，还有张材家的周瑞家的陪 着。}
}
{
    \eA{"This is really quite true," Pao-ch'ai laughed. "Whenever we've got nothing
        to do, and we talk matters over, (we're quite unanimous) that not one in a
        hundred could be picked out to equal you girls in here. The beauty is that
        each one of you possesses her own good qualities!" "In every thing, whether
        large or small, a heavenly principle rules alike," Li Wan explained.}
}
{
}

\Verse{11}
{
    \zA{又有两三个丫头在地下，倒口袋里的枣儿、倭瓜并些野菜。 众人见他进来，都 忙站起来。 刘老老因上次来过，知道平儿的身分，忙跳下地来，问：“姑娘好？”
        又说：“家里都问好。 早要来请姑奶奶的安、看姑娘来的，因为庄家忙，好容易今 年多打了两石粮食，瓜果菜蔬也丰盛，这是头一起摘下来的，并没敢卖呢，留的尖
        儿，孝敬姑奶奶、姑娘们尝尝。 姑娘们天天山珍海味的，也吃腻了，吃个野菜儿， 也算我们的穷心。” 平儿忙道：“多谢费心。”}
}
{
    \eA{"Were there, for instance, no Yüan Yang in our venerable senior's
        apartments, how would it ever do? Commencing with Madame Wang herself, who
        is it who could muster sufficient courage to expostulate with the old lady?
        Yet she plainly has the pluck to put in her remonstrances with her; and, as
        it happens, our worthy ancestor lends a patient ear to only what she says
        and no one else.}
}
{
}

\Verse{12}
{
    \zA{又让坐，自己坐了，又让：“张婶子周 大娘坐了。” 命小丫头子：“倒茶去。” 周瑞张材两家的因笑道：“姑娘今日脸上有些 春色，眼圈儿都红了。”
        平儿笑道：“可不是，我原不喝，大奶奶和姑娘们只是拉着 死灌，不得已喝了两钟，脸就红了。” 张材家的笑道：“我倒想着要喝呢，又没人让 我。
        明日再有人请姑娘，可带了我去罢。” 说着，大家都笑了。 周瑞家的道：“早起 我就看见那螃蟹了，一斤只好秤两个三个，这么两三大篓，想是有七八十斤呢。”}
}
{
    \eA{None of the others can remember what our old senior has in the way of
        clothes and head-ornaments, but she can remember everything; and, were she
        not there to look after things, there is no knowing how many would not be
        swindled away.}
}
{
}

\Verse{13}
{
    \zA{周瑞家的又道：“要是上上下下，只怕还不够！” 平儿道：“那里都吃？ 不过都是有名 儿的吃两个子。 那些散众儿的，也有摸着的，也有摸不着的。”
        刘老老道：“这些螃 蟹，今年就值五分一斤，十斤五钱，五五二两五，三五一十五，再搭上酒菜，一共 倒有二十多两银子。
        阿弥陀佛!这一顿的银子，够我们庄家人过一年了！” 平儿因问：“想是见过奶奶了？” 刘老老道：“见过了，叫我们等着呢。”}
}
{
    \eA{That child besides is so straightforward at heart, that, despite all this,
        she often puts in a good word for others, and doesn't rely upon her
        influence to look down disdainfully upon any one!" "It was only yesterday,"
        Hsi Ch'un observed with a smile, "that our dear ancestor said that she was
        ever so much better than the whole lot of us!" "She's certainly splendid!"
        P'ing Erh ventured.}
}
{
}

\Verse{14}
{
    \zA{说着， 又往窗外看天气，说道：“天好早晚了，我们也去罢，别出不去城才是饥荒呢。” 周 瑞家的道：“等着我替你瞧瞧去。”
        说着，一径去了，半日方来，笑道：“可是老老 的福来了，竟投了这两个人的缘了。” 平儿等问：“怎么样？” 周瑞家的笑道：“二
        奶奶在老太太跟前呢，我原是悄悄的告诉二奶奶：‘刘老老要家去呢，怕晚了赶不 出城去。’ 二奶奶说：‘大远的，难为他扛了些东西来，晚了就住一夜，明日再去。’
        这可不是投上二奶奶的缘了吗？}
}
{
    \eA{"How could we rise up to her standard?" "Ts'ai Hsia," Pao-yü put in, "who is
        in mother's rooms, is a good sort of girl!" "Of course she is!" T'an Ch'un
        assented. "But she's good enough as far as external appearances go, but
        inwardly she's a sly one! Madame Wang is just like a joss; she does not give
        her mind to any sort of business; but this girl is up to everything;}
}
{
}

\Verse{15}
{
    \zA{这也罢了，偏老太太又听见了，问：‘刘老老是谁？’ 二奶奶就回明白了。 老太太又说：‘我正想个积古的老人家说话儿，请了来我见见。’
        这可不是想不到的投上缘了？” 说着，催刘老老下来前去。 刘老老道：“我这生像 儿，怎么见得呢？ 好嫂子，你就说我去了罢！”
        平儿忙道：“你快去罢，不相干的。 我们老太太最是惜老怜贫的，比不得那个狂三诈四的那些人。 想是你怯上，我和周 大娘送你去。”
        说着，同周瑞家的带了刘老老往贾母这边来。}
}
{
    \eA{and it is she who in all manner of things reminds her mistress what there is
        to be done. She even knows everything, whether large or small, connected
        with Mr. Chia Cheng's staying at home or going out of doors; and when at any
        time Madame Wang forgets, she, from behind the scenes, prompts her how to
        act." "Well, never mind about her!" Li Wan suggested.}
}
{
}

\Verse{16}
{
    \zA{二门口该班的小厮们， 见了平儿出来都站起来，有两个又跑上来，赶着平儿叫“姑娘”。 平儿问道：“又说 什么？”
        那小厮笑道：“这会子也好早晚了，我妈病着，等我去请大夫。 好姑娘， 我讨半日假，可使得？” 平儿道：“你们倒好，都商量定了，一天一个，告假又不
        回奶奶，只和我胡缠。 前日住儿去了，二爷偏叫他，叫不着，我应起来了，还说我 做了情了。 你今日又来了。”
        周瑞家的道：“当真的他妈病了，姑娘也替他应着放了 他罢。”}
}
{
    \eA{"But were," she pursued, pointing at Pao-yü, "no Hsi Jen in this young
        gentleman's quarters, just you imagine what a pitch things would reach! That
        vixen Feng may truly resemble the prince Pa of the Ch'u kingdom; and she may
        have two arms strong enough to raise a tripod weighing a thousand catties,
        but had she not this maid (P'ing Erh), would she be able to accomplish
        everything so thoroughly?"}
}
{
}

\Verse{17}
{
    \zA{平儿道：“明日一早来。 听着，我还要使你呢。 再睡的日头晒着屁股再来! 你这一去，带个信儿给旺儿，就说奶奶的话，问他那剩的利钱，明日要还不交来，
        奶奶不要了，索性送他使罢。” 那小厮欢天喜地，答应去了。 平儿等来至贾母房中。 彼时大观园中姐妹们都在贾母前承奉，刘老老进去，只
        见满屋里珠围翠绕、花枝招展的，并不知都系何人。 只见一张榻上，独歪着一位老 婆婆，身后坐着一个纱罗裹的美人一般的个丫鬟在那里捶腿，凤姐儿站着正说笑。}
}
{
    \eA{"In days gone by," P'ing Erh interposed, "four servant-girls came along with
        her, but what with those who've died and those who've gone, only I remain
        like a solitary spirit." "You're, after all, the fortunate one!" Li Wan
        retorted, "but our hussey Feng too is lucky in having you! Had I not also
        once, just remember, two girls, when your senior master Chu was alive?}
}
{
}

\Verse{18}
{
    \zA{刘老老便知是贾母了，忙上来，陪着笑，拜了几拜，口里说：“请老寿星安！” 贾母 也忙欠身问好，又命周瑞家的端过椅子来坐着。 那板儿仍是怯人，不知问候。
        贾母 道：“老亲家，你今年多大年纪了？” 刘老老忙起身答道：“我今年七十五了。” 贾 母向众人道：“这么大年纪了，还这么硬朗。
        比我大好几岁呢!我要到这个年纪，还 不知怎么动不得呢。” 刘老老笑道：“我们生来是受苦的人，老太太生来是享福的。
        我们要也这么着，那些庄家活也没人做了。”}
}
{
    \eA{Am I not, you've seen for yourselves, a person to bear with people? But in
        such a surly frame of mind did I find them both day after day that, as soon
        as your senior master departed this life, I availed myself of their youth
        (to give them in marriage) and to pack both of them out of my place.}
}
{
}

\Verse{19}
{
    \zA{贾母道：“眼睛牙齿还好？” 刘老老道： “还都好，就是今年左边的槽牙活动了。” 贾母道：“我老了，都不中用了，眼也花， 耳也聋，记性也没了。
        你们这些老亲戚，我都不记得了。 亲戚们来了，我怕人笑话， 我都不会。 不过嚼的动的吃两口，睡一觉，闷了时和这些孙子孙女儿玩笑会子就完 了。”
        刘老老笑道：“这正是老太太的福了。 我们想这么着不能。” 贾母道：“什么福， 不过是老废物罢咧！” 说的大家都笑了。}
}
{
    \eA{But had either of them been good for anything and worthy to be kept, I
        would, in fact, have now had some one to give me a helping hand!" As she
        spoke, the very balls of her eyes suddenly became quite red. "Why need you
        again distress your mind?" they with one voice, exclaimed. "Isn't it better
        that we should break up?" While conversing, they rinsed their hands;}
}
{
}

\Verse{20}
{
    \zA{贾母又笑道：“我才听见凤哥儿说，你带了 好些瓜菜来，我叫他快收拾去了。 我正想个地里现结的瓜儿菜儿吃，外头买的不像 你们地里的好吃。”
        刘老老笑道：“这是野意儿，不过吃个新鲜。 依我们倒想鱼肉吃， 只是吃不起。” 贾母又道：“今日既认着了亲，别空空的就去，不嫌我这里，就住一 两天再去。
        我们也有个园子，园子里头也有果子。 你明日也尝尝，带些家去，也算 是看亲戚一趟。”}
}
{
    \eA{and, when they had agreed to go in a company to dowager lady Chia's and
        Madame Wang's and inquire after their health, the matrons and servant-maids
        swept the pavilion and collected and washed the cups and saucers. Hsi Jen
        proceeded on her way along with P'ing Erh. "Come into my room," said Hsi Jen
        to P'ing Erh, "and sit down and have another cup of tea." "I won't have any
        tea just now," P'ing Erh answered.}
}
{
}

\Verse{21}
{
    \zA{凤姐儿见贾母喜欢，也忙留道：“我们这里虽不比你们的场院大， 空屋子还有两间，你住两天，把你们那里的新闻故事儿，说些给我们老太太听听。”
        贾母笑道：“凤丫头别拿他取笑儿，他是屯里人，老实，那里搁的住你打趣？” 说 着，又命人去先抓果子给板儿吃。 板儿见人多了，又不敢吃。
        贾母又命拿些钱给他， 叫小么儿们带他外头玩去。 刘老老吃了茶，便把些乡村中所见所闻的事情说给贾母 听，贾母越发得了趣味。}
}
{
    \eA{"I'll come some other time." So saying, she was about to go off when Hsi Jen
        called out to her and stopped her. "This month's allowances," she asked,
        "haven't yet been issued, not even to our old mistress and Madame Wang; why
        is it?" Upon catching this inquiry, P'ing Erh hastily retraced her steps and
        drew near Hsi Jen.}
}
{
}

\Verse{22}
{
    \zA{正说着，凤姐儿便命人请刘老老吃晚饭，贾母又将自己的 菜拣了几样，命人送过去给刘老老吃。 凤姐知道合了贾母的心，吃了饭便又打发过来。
        鸳鸯忙命老婆子带了刘老老去 洗了澡，自己去挑了两件随常的衣裳叫给刘老老换上。 那刘老老那里见过这般行事？
        忙换了衣裳出来，坐在贾母榻前，又搜寻些话出来说。 彼时宝玉姐妹们也都在这里 坐着，他们何曾听见过这些话，自觉比那些瞽目先生说的书还好听。}
}
{
    \eA{After looking about to see that no one was in the neighbourhood, she
        rejoined in a low tone of voice, "Drop these questions at once! They're
        sure, anyhow, to be issued in a couple of days." "Why is it," smiled Hsi
        Jen, "that this gives you such a start?"}
}
{
}

\Verse{23}
{
    \zA{那刘老老虽是 个村野人，却生来的有些见识，况且年纪老了，世情上经历过的，见头一件贾母高 兴，第二件这些哥儿姐儿都爱听，便没话也编出些话来讲。
        因说道：“我们村庄上 种地种菜，每年每日，春夏秋冬，风里雨里，那里有个坐着的空儿？ 天天都是在那
        地头上做歇马凉亭，什么奇奇怪怪的事不见呢!就像旧年冬天，接连下了几天雪， 地下压了三四尺深。 我那日起的早，还没出屋门，只听外头柴草响，我想着必定有
        人偷柴草来了。}
}
{
    \eA{"This month's allowances," P'ing Erh explained to her in a whisper, "have
        long ago been obtained in advance by our mistress Secunda and given to
        people for their own purposes; and it's when the interest has been brought
        from here and there that the various sums will be lumped together and
        payment be effected. I confide this to you, but, mind, you mustn't go and
        tell any other person about it."}
}
{
}

\Verse{24}
{
    \zA{我巴着窗户眼儿一瞧，不是我们村庄上的人――”贾母道：“必定 是过路的客人们冷了，见现成的柴火抽些烤火，也是有的。” 刘老老笑道：“也并不
        是客人，所以说来奇怪。 老寿星打量什么人？ 原来是一个十七八岁极标致的个小姑 娘儿，梳着溜油儿光的头，穿着大红袄儿，白绫子裙儿。” 刚说到这里，忽听外面
        人吵嚷起来，又说：“不相干，别唬着老太太！” 贾母等听了，忙问：“怎么了？” 丫鬟回说：“南院子马棚里走了水了，不相干，已经救下去了。”}
}
{
    \eA{"Is it likely that she hasn't yet enough money for her own requirements?"
        Hsi Jen smiled. "Or is it that she's still not satisfied? And what's the use
        of her still going on bothering herself in this way?" "Isn't it so!" laughed
        P'ing Erh. "From just handling the funds for this particular item, she has,
        during these few years, so manipulated them as to turn up several hundreds
        of taels profit out of them.}
}
{
}

\Verse{25}
{
    \zA{贾母最胆小的，听 了这话，忙起身扶了人出至廊上来瞧时，只见那东南角上火光犹亮。 贾母唬得口内 念佛，又忙命人去火神跟前烧香。
        王夫人等也忙都过来请安，回说：“已经救下去 了。 老太太请进去罢。” 贾母足足的看着火光熄了，方领众人进来。
        宝玉且忙问刘老老：“那女孩儿大雪地里做什么抽柴火？ 倘或冻出病来呢？” 贾 母道：“都是才说抽柴火，惹出事来了，你还问呢!别说这个了，说别的罢。” 宝玉
        听说，心内虽不乐，也只得罢了。}
}
{
    \eA{Nor does she spend that monthly allowance of hers for public expenses. But
        the moment she accumulates anything like eight or ten taels odd, she gives
        them out too. Thus the interest on her own money alone comes up to nearly a
        thousand taels a year." "You and your mistress take our money," Hsi Jen
        observed laughingly, "and get interest on it; fooling us as if we were no
        better than idiots."}
}
{
}

\Verse{26}
{
    \zA{刘老老便又想了想，说道：“我们庄子东边庄上 有个老奶奶子，今年九十多岁了。 他天天吃斋念佛，谁知就感动了观音菩萨，夜里
        来托梦，说：‘你这么虔心，原本你该绝后的，如今奏了玉皇，给你个孙子。’ 原来 这老奶奶只有一个儿子，这儿子也只一个儿子，好容易养到十七八岁上，死了，哭
        的什么儿似的。 后起间，真又养了一个，今年才十三四岁，长得粉团儿似的，聪明 伶俐的了不得呢。 这些神佛是有的不是！”}
}
{
    \eA{"Here you are again with your uncharitable words!" P'ing Erh remonstrated.
        "Can it be that you haven't yet enough to meet your own expenses with?" "I
        am, it's true, not short of money," Hsi Jen replied, "as I have nowhere to
        go and spend it; but the thing is that I'm making provision for that fellow
        of ours, (Pao-yü)."}
}
{
}

\Verse{27}
{
    \zA{这一席话暗合了贾母王夫人的心事，连 王夫人也都听住了。 宝玉心中只惦记抽柴的事，因闷的心中筹画。 探春因问他：“昨日扰了史大妹
        妹，咱们回去商议着邀一社，又还了席，也请老太太赏菊何如？” 宝玉笑道：“老 太太说了，还要摆酒还史妹妹的席，叫咱们做陪呢。 等吃了老太太的，咱们再请不
        迟。” 探春道：“越往前越冷了，老太太未必高兴。” 宝玉道：“老太太又喜欢下雨下 雪的，咱们等下头场雪，请老太太赏雪不好吗？
        咱们雪下吟诗，也更有趣了。”}
}
{
    \eA{"If you ever find yourself in any great straits and need money," P'ing Erh
        resumed, "you're at liberty to take first those few taels I've got over
        there to suit your own convenience with, and by and bye I can reduce them
        from what is due to you and we'll be square." "I'm not in need of any just
        now," retorted Hsi Jen. "But should I not have enough, when I want some,
        I'll send some one to fetch them, and finish."}
}
{
}

\Verse{28}
{
    \zA{黛玉 笑道：“咱们雪下吟诗，依我说，还不如弄一捆柴火，雪下抽柴，还更有趣儿呢！” 说着，宝钗等都笑了。 宝玉瞅了他一眼，也不答话。
        一时散了，背地里宝玉到底拉了刘老老，细问那女孩儿是谁。 刘老老只得编了 告诉他：“那原是我们庄子北沿儿地埂子上，有个小祠堂儿，供的不是神佛，当先
        有个什么老爷――”说着，又想名姓。 宝玉道：“不拘什么名姓，也不必想了，只 说原故就是了。”}
}
{
    \eA{P'ing Erh promised that she would let her have the money at any time she
        sent for it, and, and taking the shortest cut, she issued out of the garden
        gate. Here she encountered a servant despatched from the other side by lady
        Feng. She came in search of P'ing Erh. "Our lady," she said, "has something
        for you to do, and is waiting for you." "What's up that it's so pressing?"
        P'ing Erh inquired.}
}
{
}

\Verse{29}
{
    \zA{刘老老道：“这老爷没有儿子，只有一位小姐，名字叫什么若玉， 知书儿识字的，老爷太太爱的像珍珠儿。 可惜了儿的，这小姐儿长到十七岁了，一 病就病死了。”
        宝玉听了，跌足叹惜，又问：“后来怎么样？” 刘老老道：“因为老 爷太太疼的心肝儿似的，盖了那祠堂，塑了个像儿，派了人烧香儿拨火的。 如今年
        深日久了，人也没了，庙也烂了，那泥胎儿可就成了精咧。” 宝玉忙道：“不是成精， 规矩这样人是不死的。” 刘老老道：“阿弥陀佛!是这么着吗？}
}
{
    \eA{"Our senior mistress detained me by force to have a chat, so I couldn't
        manage to get away. But here she time after time sends people after me in
        this manner!" "Whether you go or not is your own look out," the maid
        replied. "It isn't worth your while getting angry with me! If you dare, go
        and tell these things to our mistress!" P'ing Erh spat at her
        contemptuously, and rushed back in anxious haste.}
}
{
}

\Verse{30}
{
    \zA{不是哥儿说，我们还当 他成了精了呢。 他时常变了人出来闲逛。 我才说抽柴火的，就是他了。 我们村庄上 的人商量着还要拿榔头砸他呢。”
        宝玉忙道：“快别如此。 要平了庙，罪过不小！” 刘老老道：“幸亏哥儿告诉我，明日回去，拦住他们就是了。” 宝玉道：“我们老太
        太、太太都是善人，就是合家大小也都好善喜舍，最爱修庙塑神的。 我明日做一个 疏头，替你化些布施，你就做香头，攒了钱，把这庙修盖，再装塑了泥像，每月给
        你香火钱烧香，好不好？”}
}
{
    \eA{She discovered, however, that lady Feng was not at home. But unexpectedly
        she perceived that the old goody Liu, who had paid them a visit on a
        previous occasion for the purpose of obtaining pecuniary assistance, had
        come again with Pan Erh, and was seated in the opposite room, along with
        Chang Ts'ai's wife and Chou Jui's wife, who kept her company.}
}
{
}

\Verse{31}
{
    \zA{刘老老道：“若这样时，我托那小姐的福，也有几个钱 使了。” 宝玉又问他地名庄名，来往远近，坐落何方，刘老老便顺口诌了出来。
        宝玉信以为真，回至房中，盘算了一夜。 次日一早，便出来给了焙茗几百钱， 按着刘老老说的方向地名，着焙茗去先踏看明白，回来再作主意。 那焙茗去后，宝
        玉左等也不来，右等也不来，急的热地里的蚰蜒似的。 好容易等到日落，方见焙茗 兴兴头头的回来了。 宝玉忙问：“可找着了？”}
}
{
    \eA{But two or three servant-maids were inside as well emptying on the floor
        bags containing dates, squash and various wild greens. As soon as they saw
        her appear in the room, they promptly stood up in a body. Old goody Liu had,
        on her last visit, learnt what P'ing Erh's status in the establishment was,
        so vehemently jumping down, she enquired, "Miss, how do you do? All at
        home," she pursued, "send you their compliments.}
}
{
}

\Verse{32}
{
    \zA{焙茗笑道：“爷听的不明白，叫我好 找!那地名坐落，不像爷听的一样，所以找了一天，找到东北角田埂子上，才有一 个破庙。”
        宝玉听说，喜的眉开眼笑，忙说道：“刘老老有年纪的人，一时错记了也 是有的。 你且说你见的。” 焙茗道：“那庙门却倒也朝南开，也是稀破的。 我找的正
        没好气，一见这个，我说可好了，连忙进去。 一看泥胎，唬的我又跑出来了，活像 真的似的！” 宝玉喜的笑道：“他能变化人了，自然有些生气。”
        焙茗拍手道：“那里 是什么女孩儿？}
}
{
    \eA{I meant to have come earlier and paid my respects to my lady and to look you
        up, miss; but we've been very busy on the farm. We managed this year to
        reap, after great labour, a few more piculs of grain than usual. But melons,
        fruits and vegetables have also been plentiful. These things, you see here,
        are what we picked during the first crop;}
}
{
}

\Verse{33}
{
    \zA{竟是一位青脸红发的瘟神爷！” 宝玉听了，啐了一口，骂道：“真是 个没用的杀材，这点子事也干不来！” 焙茗道：“爷又不知看了什么书，或者听了谁
        的混帐语，信真了，把这件没头脑的事派我去碰头。 怎么说我没用呢？” 宝玉见他 急了，忙抚慰他道：“你别急，改日闲了，你再找去。
        要是他哄我们呢，自然没了； 要竟是有的，你岂不也积了阴骘呢？ 我必重重的赏你。” 说着，只见二门上的小厮来
        说：“老太太屋里的姑娘们站在二门口找二爷呢。”}
}
{
    \eA{and as we didn't presume to sell them, we kept the best to present to our
        lady and the young ladies to taste. The young ladies must, of course, be
        surfeited with all the delicacies and fine things they daily get, but by
        having some of our wild greens to eat, they will show some regard for our
        poor attention." "Many thanks for all the trouble you have taken!" Ping Erh
        eagerly rejoined.}
}
{
}

\Verse{34}
{
    \zA{不知何事，下回分解。 (本章完)}
}
{
    \eA{Then pressing her to resume her place, she sat down herself; and, urging
        Mrs. Chang and Mrs. Chou to take their seats, she bade a young waiting-maid
        go and serve the tea. "There's a joyous air about your face to-day, Miss,
        and your eye-balls are all red," the wife of Chou Jui and the wife of Chang
        Ts'ai thereupon smilingly ventured. "Naturally!" P'ing Erh laughed.}
}
{
}

\Verse{35}
{
    \zA{}
}
{
    \eA{"I generally don't take any wine, but our senior mistress, and our young
        ladies caught hold of me and insisted upon pouring it down my throat. I had
        no alternative therefore but to swallow two cups full; so my face at once
        flushed crimson." "I have a longing for wine," Chang Ts'ai's wife smiled;
        "but there's no one to offer me any. But when any one by and by invites you,
        Miss, do take me along with you!" At these words, one and all burst out
        laughing. "Early this morning," Chou Jui's wife interposed, "I caught a
        glimpse of those crabs. Only two or three of them would weigh a catty; so in
        those two or three huge hampers, there must have been, I presume, seventy to
        eighty catties!" "If some were intended for those above as well as for those
        below;"}
}
{
}

\Verse{36}
{
    \zA{}
}
{
    \eA{Chou Jui's wife added, "they couldn't, nevertheless, I fear, have been
        enough." "How could every one have had any?" P'ing Erh observed. "Those
        simply with any name may have tasted a couple of them; but, as for the rest,
        some may have touched them with the tips of their hands, but many may even
        not have done as much." "Crabs of this kind!" put in old goody Liu, "cost
        this year five candareens a catty; ten catties for five mace; five times
        five make two taels five, and three times five make fifteen; and adding what
        was wanted for wines and eatables, the total must have come to something
        over twenty taels. O-mi-to-fu! why, this heap of money is ample for us
        country-people to live on through a whole year!" "I expect you have seen our
        lady?" P'ing Erh then asked. "Yes, I have seen her," assented old goody Liu.}
}
{
}

\Verse{37}
{
    \zA{}
}
{
    \eA{"She bade us wait." As she spoke, she again looked out of the window to see
        what the time of the day could be. "It's getting quite late," she afterwards
        proceeded. "We must be going, or else we mayn't be in time to get out of the
        city gates; and then we'll be in a nice fix." "Quite right," Chou Jui's wife
        observed. "I'll go and see what she's up to for you." With these words, she
        straightway left the room. After a long absence, she returned. "Good fortune
        has, indeed, descended upon you, old dame!" she smiled. "Why, you've won the
        consideration of those two ladies!" "What about it?" laughingly inquired
        P'ing Erh and the others. "Lady Secunda," Chou Jui's wife explained with a
        smile, "was with our venerable lady, so I gently whispered to her: 'old
        goody Liu wishes to go home;}
}
{
}

\Verse{38}
{
    \zA{}
}
{
    \eA{it's getting late and she fears she mightn't be in time to go out of the
        gates!' 'It's such a long way off!' Our lady Secunda rejoined, 'and she had
        all the trouble and fatigue of carrying that load of things; so if it's too
        late, why, let her spend the night here and start on the morrow!' Now isn't
        this having enlisted our mistress' sympathies? But not to speak of this! Our
        old lady also happened to overhear what we said, and she inquired: 'who is
        old goody Liu?' Our lady Secunda forthwith told her all. 'I was just
        longing,' her venerable ladyship pursued, 'for some one well up in years to
        have a chat with; ask her in, and let me see her!' So isn't this coming in
        for consideration, when least unexpected?"}
}
{
}

\Verse{39}
{
    \zA{}
}
{
    \eA{So speaking, she went on to urge old goody Liu to get down and betake
        herself to the front. "With a figure like this of mine," old goody Liu
        demurred, "how could I very well appear before her? My dear sister-in-law,
        do tell her that I've gone!" "Get on! Be quick!" P'ing Erh speedily cried.
        "What does it matter? Our old lady has the highest regard for old people and
        the greatest pity for the needy! She's not one you could compare with those
        haughty and overbearing people! But I fancy you're a little too timid, so
        I'll accompany you as far as there, along with Mrs. Chou." While tendering
        her services, she and Chou Jui's wife led off old goody Liu and crossed over
        to dowager lady Chia's apartments on this side of the mansion. The
        boy-servants on duty at the second gate stood up when they saw P'ing Erh
        approach.}
}
{
}

\Verse{40}
{
    \zA{}
}
{
    \eA{But two of them also ran up to her, and, keeping close to her heels: "Miss!"
        they shouted out. "Miss!" "What have you again got to say?" P'ing Erh asked.
        "It's pretty late just now," one of the boys smilingly remarked; "and mother
        is ill and wants me to go and call the doctor, so I would, dear Miss, like
        to have half a day's leave; may I?" "Your doings are really fine!" P'ing Erh
        exclaimed. "You've agreed among yourselves that each day one of you should
        apply for furlough; but instead of speaking to your lady, you come and
        bother me! The other day that Chu Erh went, Mr. Secundus happened not to
        want him, so I assented, though I also added that I was doing it as a
        favour; but here you too come to-day!" "It's quite true that his mother is
        sick," Chou Jui's wife interceded;}
}
{
}

\Verse{41}
{
    \zA{}
}
{
    \eA{"so, Miss, do say yes to him also, and let him go!" "Be back as soon as it
        dawns to-morrow!" P'ing Erh enjoined. "Wait, I've got something for you to
        do, for you'll again sleep away, and only turn up after the sun has blazed
        away on your buttocks. As you go now, give a message to Wang Erh! Tell him
        that our lady bade you warn him that if he does not hand over the balance of
        the interest due by to-morrow, she won't have anything to do with him. So
        he'd better let her have it to meet her requirements and finish." The
        servant-lad felt in high glee and exuberant spirits. Expressing his
        obedience, he walked off. P'ing Erh and her companions repaired then to old
        lady Chia's apartments.}
}
{
}

\Verse{42}
{
    \zA{}
}
{
    \eA{Here the various young ladies from the Garden of Broad Vista were at the
        time assembled paying their respects to their grandmother. As soon as old
        goody Liu put her foot inside, she saw the room thronged with girls (as
        seductive) as twigs of flowers waving to and fro, and so richly dressed, as
        to look enveloped in pearls, and encircled with king-fisher ornaments. But
        she could not make out who they all were. Her gaze was, however, attracted
        by an old dame, reclining alone on a divan. Behind her sat a girl, a regular
        beauty, clothed in gauze, engaged in patting her legs. Lady Feng was on her
        feet in the act of cracking some joke. Old goody Liu readily concluded that
        it must be dowager lady Chia, so promptly pressing forward, she put on a
        forced smile and made several curtseys. "My obeisance to you, star of
        longevity!"}
}
{
}

\Verse{43}
{
    \zA{}
}
{
    \eA{she said. Old lady Chia hastened, on her part, to bow and to inquire after
        her health. Then she asked Chou Jui's wife to bring a chair over for her to
        take a seat. But Pan Erh was still so very shy that he did not know how to
        make his obeisance. "Venerable relative," dowager lady Chia asked, "how old
        are you this year?" Old goody Liu immediately rose to her feet. "I'm
        seventy-five this year," she rejoined. "So old and yet so hardy!" Old lady
        Chia remarked, addressing herself to the party. "Why she's older than myself
        by several years! When I reach that age, I wonder whether I shall be able to
        move!" "We people have," old goody Liu smilingly resumed, "to put up, from
        the moment we come into the world, with ever so many hardships; while your
        venerable ladyship enjoys, from your birth, every kind of blessing!}
}
{
}

\Verse{44}
{
    \zA{}
}
{
    \eA{Were we also like this, there'd be no one to carry on that farming work."
        "Are your eyes and teeth still good?" Dowager lady Chia went on to inquire.
        "They're both still all right," old goody Liu replied. "The left molars,
        however, have got rather shaky this year." "As for me, I'm quite an old
        fossil," dowager lady Chia observed. "I'm no good whatever. My eyesight is
        dim; my ears are deaf, my memory is gone. I can't even recollect any of you,
        old family connections. When therefore any of our relations come on a visit,
        I don't see them for fear lest I should be ridiculed. All I can manage to
        eat are a few mouthfuls of anything tender enough for my teeth; and I can
        just dose a bit or, when I feel in low spirits, I distract myself a little
        with these grandsons and grand-daughters of mine; that's all I'm good for."}
}
{
}

\Verse{45}
{
    \zA{}
}
{
    \eA{"This is indeed your venerable ladyship's good fortune!" old goody Liu
        smiled. "We couldn't enjoy anything of the kind, much though we may long for
        it." "What good fortune!" dowager lady Chia exclaimed. "I'm a useless old
        thing, no more." This remark made every one explode into laughter. Dowager
        lady Chia also laughed. "I heard our lady Feng say a little while back," she
        added, "that you had brought a lot of squash and vegetables, and I told her
        to put them by at once. I had just been craving to have newly-grown melons
        and vegetables; but those one buys outside are not as luscious as those
        produced in your farms." "This is the rustic notion," old goody Liu laughed,
        "to entirely subsist on fresh things!}
}
{
}

\Verse{46}
{
    \zA{}
}
{
    \eA{Yet, we long to have fish and meat for our fare, but we can't afford it."
        "I've found a relative in you to-day," dowager lady Chia said, "so you
        shouldn't go empty-handed! If you don't despise this place as too mean, do
        stay a day or two before you start! We've also got a garden here; and this
        garden produces fruits too; you can taste some of them to-morrow and take a
        few along with you home, in order to make it look like a visit to
        relatives." When lady Feng saw how delighted old lady Chia was with the
        prospects of the old dame's stay, she too lost no time in doing all she
        could to induce her to remain. "Our place here," she urged, "isn't, it's
        true, as spacious as your threshing-floor;}
}
{
}

\Verse{47}
{
    \zA{}
}
{
    \eA{but as we've got two vacant rooms, you'd better put up in them for a couple
        of days, and choose some of your village news and old stories and recount
        them to our worthy senior." "Now you, vixen Feng," smiled dowager lady Chia,
        "don't raise a laugh at her expense! She's only a country woman; and will an
        old dame like her stand any chaff from you?" While remonstrating with her,
        she bade a servant go, before attending to anything else, and pluck a few
        fruits. These she handed to Pan Erh to eat. But Pan Erh did not venture to
        touch them, conscious as he was of the presence of such a number of
        bystanders. So old lady Chia gave orders that a few cash should be given
        him, and then directed the pages to take him outside to play.}
}
{
}

\Verse{48}
{
    \zA{}
}
{
    \eA{After sipping a cup of tea, old goody Liu began to relate, for the benefit
        of dowager lady Chia, a few of the occurrences she had seen or heard of in
        the country. These had the effect of putting old lady Chia in a more
        exuberant frame of mind. But in the midst of her narration, a servant, at
        lady Feng's instance, asked goody Liu to go and have her evening meal.
        Dowager lady Chia then picked out, as well, several kinds of eatables from
        her own repast, and charged some one to take them to goody Liu to feast on.
        But the consciousness that the old dame had taken her senior's fancy induced
        lady Feng to send her back again as soon as she had taken some refreshments.
        On her arrival, Yüan Yang hastily deputed a matron to take goody Liu to have
        a bath.}
}
{
}

\Verse{49}
{
    \zA{}
}
{
    \eA{She herself then went and selected two pieces of ordinary clothes, and these
        she entrusted to a servant to hand to the old dame to change. Goody Liu had
        hitherto not set eyes upon any such grand things, so with eagerness she
        effected the necessary alterations in her costume. This over, she made her
        appearance outside, and, sitting in front of the divan occupied by dowager
        lady Chia, she went on to narrate as many stories as she could recall to
        mind.}
}
{
}

\Verse{50}
{
    \zA{}
}
{
    \eA{Pao-yü and his cousins too were, at the time, assembled in the room, and as
        they had never before heard anything the like of what she said, they, of
        course, thought her tales more full of zest than those related by itinerant
        blind story-tellers. Old goody Liu was, albeit a rustic person, gifted by
        nature with a good deal of discrimination. She was besides advanced in
        years; and had gone through many experiences in her lifetime, so when she,
        in the first place, saw how extremely delighted old lady Chia was with her,
        and, in the second, how eager the whole crowd of young lads and lasses were
        to listen to what fell from her mouth, she even invented, when she found her
        own stock exhausted, a good many yarns to recount to them.}
}
{
}

\Verse{51}
{
    \zA{}
}
{
    \eA{"What with all the sowing we have to do in our fields and the vegetables we
        have to plant," she consequently proceeded, "have we ever in our village any
        leisure to sit with lazy hands from year to year and day to day; no matter
        whether it's spring, summer, autumn or winter, whether it blows or whether
        it rains? Yea, day after day all that we can do is to turn the bare road
        into a kind of pavilion to rest and cool ourselves on! But what strange
        things don't we see! Last winter, for instance, snow fell for several
        consecutive days, and it piled up on the ground three or four feet deep. One
        day, I got up early, but I hadn't as yet gone out of the door of our house
        when I heard outside the noise of firewood (being moved).}
}
{
}

\Verse{52}
{
    \zA{}
}
{
    \eA{I fancied that some one must have come to steal it, so I crept up to a hole
        in the window; but, lo, I discovered that it was no one from our own
        village." "It must have been," interposed dowager lady Chia, "some
        wayfarers, who being smitten with the cold, took some of the firewood, they
        saw ready at hand, to go and make a fire and warm themselves with! That's
        highly probable!" "It was no wayfarers at all," old goody Liu retorted
        smiling, "and that's what makes the story so strange. Who do you think it
        was, venerable star of longevity? It was really a most handsome girl of
        seventeen or eighteen, whose hair was combed as smooth as if oil had been
        poured over it. She was dressed in a deep red jacket, a white silk
        petticoat…."}
}
{
}

\Verse{53}
{
    \zA{}
}
{
    \eA{When she reached this part of her narrative, suddenly became audible the
        voices of people bawling outside. "It's nothing much," they shouted, "don't
        frighten our old mistress!" Dowager lady Chia and the other inmates caught,
        however, their cries and hurriedly inquired what had happened. A
        servant-maid explained in reply that a fire had broken out in the stables in
        the southern court, but that there was no danger, as the flames had been
        suppressed. Their old grandmother was a person with very little nerve. The
        moment, therefore, the report fell on her car, she jumped up with all
        despatch, and leaning on one of the family, she rushed on to the verandah to
        ascertain the state of things.}
}
{
}

\Verse{54}
{
    \zA{}
}
{
    \eA{At the sight of the still brilliant light, shed by the flames, on the south
        east part of the compound, old lady Chia was plunged in consternation, and
        invoking Buddha, she went on to shout to the servants to go and burn incense
        before the god of fire. Madame Wang and the rest of the members of the
        household lost no time in crossing over in a body to see how she was getting
        on. "The fire has been already extinguished," they too assured her, "please,
        dear ancestor, repair into your rooms!" But it was only after old lady Chia
        had seen the light of the flames entirely subside that she at length led the
        whole company indoors. "What was that girl up to, taking the firewood in
        that heavy fall of snow?" Pao-yü thereupon vehemently inquired of goody Liu.}
}
{
}

\Verse{55}
{
    \zA{}
}
{
    \eA{"What, if she had got frostbitten and fallen ill?" "It was the reference
        made recently to the firewood that was being abstracted," his grandmother
        Chia said, "that brought about this fire; and do you still go on asking more
        about it? Leave this story alone, and tell us something else!" Hearing this
        reminder, Pao-yü felt constrained to drop the subject, much against his
        wishes, and old goody Liu forthwith thought of something else to tell them.
        "In our village," she resumed, "and on the eastern side of our farmstead,
        there lives an old dame, whose age is this year, over ninety. She goes in
        daily for fasting, and worshipping Buddha.}
}
{
}

\Verse{56}
{
    \zA{}
}
{
    \eA{Who'd have thought it, she so moved the pity of the goddess of mercy that
        she gave her this message in a dream: 'It was at one time ordained that you
        should have no posterity, but as you have proved so devout, I have now
        memorialised the Pearly Emperor to grant you a grandson!' The fact is, this
        old dame had one son. This son had had too an only son; but he died after
        they had with great difficulty managed to rear him to the age of seventeen
        or eighteen. And what tears didn't they shed for him! But, in course of
        time, another son was actually born to him. He is this year just thirteen or
        fourteen, resembles a very ball of flower, (so plump is he), and is clever
        and sharp to an exceptional degree!}
}
{
}

\Verse{57}
{
    \zA{}
}
{
    \eA{So this is indeed a clear proof that those spirits and gods do exist!" This
        long tirade proved to be in harmony with dowager lady Chia's and Madame
        Wang's secret convictions on the subject. Even Madame Wang therefore
        listened to every word with all profound attention. Pao-yü, however, was so
        pre-occupied with the story about the stolen firewood that he fell in a
        brown study and gave way to conjectures. "Yesterday," T'an Ch'un at this
        point remarked, "We put cousin Shih to a lot of trouble and inconvenience,
        so, when we get back, we must consult about convening a meeting, and, while
        returning her entertainment, we can also invite our venerable ancestor to
        come and admire the chrysanthemums; what do you think of this?"}
}
{
}

\Verse{58}
{
    \zA{}
}
{
    \eA{"Our worthy senior," smiled Pao-yü, "has intimated that she means to give a
        banquet to return cousin Shih's hospitality, and to ask us to do the
        honours. Let's wait therefore until we partake of grandmother's collation,
        before we issue our own invitations; there will be ample time then to do
        so." "The later it gets, the cooler the weather becomes," T'an Ch'un
        observed, "and our dear senior is not likely to enjoy herself."
        "Grandmother," added Pao-yü, "is also fond of rain and snow, so wouldn't it
        be as well to wait until the first fall, and then ask her to come and look
        at the snow. This will be better, won't it? And were we to recite our verses
        with snow about us, it will be ever so much more fun!"}
}
{
}

\Verse{59}
{
    \zA{}
}
{
    \eA{"To hum verses in the snow," Lin Tai-yü speedily demurred with a smile,
        "won't, in my idea, be half as nice as building up a heap of firewood and
        then stealing it, with the flakes playing about us. This will be by far more
        enjoyable!" This proposal made Pao-ch'ai and the others laugh. Pao-yü cast a
        glance at her but made no reply. But, in a short time, the company broke up.
        Pao-yü eventually gave old goody Liu a tug on the sly and plied her with
        minute questions as to who the girl was. The old dame was placed under the
        necessity of fabricating something for his benefit. "The truth is," she
        said, "that there stands on the north bank of the ditch in our village a
        small ancestral hall, in which offerings are made, but not to spirits or
        gods.}
}
{
}

\Verse{60}
{
    \zA{}
}
{
    \eA{There was in former days some official or other…" "While speaking, she went
        on to try and recollect his name and surname. "No matter about names or
        surnames!" Pao-yü expostulated. "There's no need for you to recall them to
        memory! Just mention the facts; they'll be enough." "This official," old
        goody Liu resumed, "had no son. His offspring consisted of one young
        daughter, who went under the name of Jo Yü, (like Jade). She could read and
        write, and was doated upon by this official and his consort, just as if she
        were a precious jewel. But, unfortunately, when this young lady, Jo Yü, grew
        up to be seventeen, she contracted some disease and died." When these words
        fell on Pao-yü's ears, he stamped his foot and heaved a sigh.}
}
{
}

\Verse{61}
{
    \zA{}
}
{
    \eA{"What happened after that?" he then asked. Old goody Liu pursued her story.
        "So incessantly," she continued, "did this official and his consort think of
        their child that they raised this ancestral hall, erected a clay image of
        their young daughter Jo Yü in it, and appointed some one to burn incense and
        trim the fires. But so many days and years have now elapsed that the people
        themselves are no more alive, the temple is in decay, and the image itself
        is become a spirit." "It hasn't become a spirit," remonstrated Pao-yü with
        vehemence. "Human beings of this kind may, the rule is, die, yet they are
        not dead." "O-mi-to-fu!" ejaculated old goody Liu; "is it really so!}
}
{
}

\Verse{62}
{
    \zA{}
}
{
    \eA{Had you, sir, not enlightened us, we would have remained under the
        impression that she had become a spirit! But she repeatedly transforms
        herself into a human being, and there she roams about in every village,
        farmstead, inn and roadside. And the one I mentioned just now as having
        taken the firewood is that very girl! The villagers in our place are still
        consulting with the idea of breaking this clay image and razing the temple
        to the ground." "Be quick and dissuade them!" eagerly exclaimed Pao-yü.
        "Were they to raze the temple to the ground, their crime won't be small."
        "It's lucky that you told me, Sir," old goody Liu added. "When I get back
        to-morrow, I'll make them relinquish the idea and finish!"}
}
{
}

\Verse{63}
{
    \zA{}
}
{
    \eA{"Our venerable senior and my mother," Pao-yü pursued, "are both charitable
        persons. In fact, all the inmates of our family, whether old or young, do,
        in like manner, delight in good deeds, and take pleasure in distributing
        alms. Their greatest relish is to repair temples, and to put up images to
        the spirits; so to-morrow, I'll make a subscription and collect a few
        donations for you, and you can then act as incense-burner. When sufficient
        money has been raised, this fane can be repaired, and another clay image put
        up; and month by month I'll give you incense and fire money to enable you to
        burn joss-sticks; won't this be A good thing for you?"}
}
{
}

\Verse{64}
{
    \zA{}
}
{
    \eA{"In that case," old goody Liu rejoined, "I shall, thanks to that young
        lady's good fortune, have also a few cash to spend." Pao-yü thereupon
        likewise wanted to know what the name of the place was, the name of the
        village, how far it was there and back, and whereabout the temple was
        situated. Old goody Liu replied to his questions, by telling him every idle
        thought that came first to her lips. Pao-yü, however, credited the
        information she gave him and, on his return to his rooms, he exercised, the
        whole night, his mind with building castles in the air.}
}
{
}

\Verse{65}
{
    \zA{}
}
{
    \eA{On the morrow, as soon as daylight dawned, he speedily stepped out of his
        room, and, handing Pei Ming several hundreds of cash, he bade him proceed
        first in the direction and to the place specified by old goody Liu, and
        clearly ascertain every detail, so as to enable him, on his return from his
        errand, to arrive at a suitable decision to carry out his purpose. After Pei
        Ming's departure, Pao-yü continued on pins on needles and on the tiptoe of
        expectation. Into such a pitch of excitement did he work himself, that he
        felt like an ant in a burning pan. With suppressed impatience, he waited and
        waited until sunset. At last then he perceived Pei Ming walk in, in high
        glee. "Have you discovered the place?" hastily inquired Pao-yü.}
}
{
}

\Verse{66}
{
    \zA{}
}
{
    \eA{"Master," Pei Ming laughed, "you didn't catch distinctly the directions
        given you, and you made me search in a nice way! The name of the place and
        the bearings can't be those you gave me, Sir; that is why I've had to hunt
        about the whole day long! I prosecuted my inquiries up to the very ditch on
        the north east side, before I eventually found a ruined temple." Upon
        hearing the result of his researches, Pao-yü was much gratified. His very
        eyebrows distended. His eyes laughed. "Old goody Liu," he said with
        eagerness, "is a person well up in years, and she may at the moment have
        remembered wrong; it's very likely she did. But recount to me what you saw."
        "The door of that temple," Pei Ming explained, "really faces south, and is
        all in a tumble-down condition.}
}
{
}

\Verse{67}
{
    \zA{}
}
{
    \eA{I searched and searched till I was driven to utter despair. As soon,
        however, as I caught sight of it, 'that's right,' I shouted, and promptly
        walked in. But I at once discovered a clay figure, which gave me such a
        fearful start, that I scampered out again; for it looked as much alive as if
        it were a real living being." Pao-yü smiled full of joy. "It can
        metamorphose itself into a human being," he observed, "so, of course, it has
        more or less a life-like appearance." "Was it ever a girl?" Pei Ming
        rejoined clapping his hands. "Why it was, in fact, no more than a
        green-faced and red-haired god of plagues." Pao-yü, at this answer, spat at
        him contemptuously. "You are, in very truth, a useless fool!" he cried.}
}
{
}

\Verse{68}
{
    \zA{}
}
{
    \eA{"Haven't you even enough gumption for such a trifling job as this?" "What
        book, I wonder, have you again been reading, master?" Pei Ming continued.
        "Or you may, perhaps, have heard some one prattle a lot of trash and
        believed it as true! You send me on this sort of wild goose chase and make
        me go and knock my head about, and how can you ever say that I'm good for
        nothing?" Pao-yü did not fail to notice that he was in a state of
        exasperation so he lost no time in trying to calm him. "Don't be impatient!"
        he urged. "You can go again some other day, when you've got nothing to
        attend to, and institute further inquiries! If it turns out that she has
        hood-winked us, why, there will, naturally, be no such thing.}
}
{
}

\Verse{69}
{
    \zA{}
}
{
    \eA{But if, verily, there is, won't you also lay up for yourself a store of good
        deeds? I shall feel it my duty to reward you in a most handsome manner." As
        he spoke, he espied a servant-lad, on service at the second gate, approach
        and report to him: "The young ladies in our venerable ladyship's apartments
        are standing at the threshold of the second gate and looking out for you,
        Mr. Secundus." But as, reader, you are not aware what they were on the
        look-out to tell him, the subsequent chapter will explain it for you.}
}
{
}
